\documentclass[10pt,letterpaper]{article}
\usepackage{cornell-notes}

\title{Chapter 6: Random Composition}
\author{}
\date{}

\setDocTitle{Chapter 6: Random Composition}
\setDocSubtitle{Combinatorial Algorithms}
\setDocVersion{v1.0}
\setDocStatus{Study Companion}
\setDocOwner{Combinatorial Algorithms Cornell Notes}
\setDocAuthor{}
\setDocDate{\today}
\setDocSummary{Cornell-format study and review notes for Chapter 6: Random Composition.}

\setCornellCollection{Combinatorial Algorithms Cornell Notes}
\setCornellUnitType{Chapter}
\setCornellUnitNumber{6}
\setCornellUnitTitle{Random Composition}
\setCornellCompanionLabel{Study and review companion}

\hypersetup{pdftitle={Chapter 6: Random Composition},pdfsubject={Cornell notes},pdfauthor={}}

\begin{document}
\maketitle
\begin{tcolorbox}[colback=Paper,colframe=Ink]{\small\bfseries\color{Teal}CORNELL NOTES}\par{\LARGE\bfseries Chapter 6: Random Composition of $n$ into $k$ Parts (RANCOM)}\par\medskip\textbf{Topic:} Uniform weak-composition sampling\hfill\textbf{Date:} \today\par\textbf{Study purpose:} Convert a uniform separator subset into a uniform composition.\end{tcolorbox}
\begin{tcolorbox}[colback=Teal!5,colframe=Teal,title=Essential question]Why does uniformly selecting separator locations and differencing them produce a uniform random weak composition?\end{tcolorbox}
\begin{tcolorbox}[colback=white,colframe=Mist,title=Learning objectives]\begin{itemize}\item Apply the balls-and-bars bijection probabilistically.\item Convert sorted separator positions into nonnegative parts.\item State the uniformity argument.\item Relate RANCOM to RANKSB.\item Handle the endpoint sentinels correctly.\end{itemize}\textbf{Prerequisites:} weak compositions, combinations, uniform sampling without replacement.\end{tcolorbox}
\section*{Cornell notes}
\begin{longtable}{@{}Q!{\color{Mist}\vrule width 1pt}A@{}}\cornellhead
\cue{What is sampled?}{Select a uniformly random $(k-1)$-subset $a_1<\cdots<a_{k-1}$ of $\{1,\ldots,n+k-1\}$. These are the interior separator locations in the balls-and-bars representation.}
\cue{How are parts recovered?}{Insert conceptual sentinels $a_0=0$ and $a_k=n+k$. Then
\[
r_i=a_i-a_{i-1}-1\qquad(1\le i\le k).
\]
Each gap counts the balls between consecutive separators. The telescope gives $\sum_i r_i=n$, and strict separator order gives $r_i\ge0$.}
\cue{Why is the output uniform?}{The separator map is a bijection between $(k-1)$-subsets of $n+k-1$ positions and weak compositions of $n$ into $k$ parts. A uniform sample in the domain of a bijection maps to a uniform sample in its codomain. Every composition therefore has probability $\binom{n+k-1}{k-1}^{-1}$.}
\cue{What does RANCOM call?}{RANCOM delegates the subset choice to RANKSB with universe size $n+k-1$ and sample size $k-1$. It then reuses the returned separator array, appends the final sentinel, and replaces the entries with successive differences minus one.}
\cue{Cost and storage}{The source states average $O(k)$ work per composition, inherited from the subset selection in the intended regime plus one differencing pass. The output array of length $k$ is also the working array.}
\cue{Boundary cases}{For $k=1$, no interior separator is chosen and the only composition is $(n)$. For $n=0$, all $k$ parts are zero if the surrounding implementation admits zero. The sentinel must be $n+k$, not $n+k-1$, because one is subtracted from every gap.}
\cue{Historical or legacy context}{The routine relies on RANKSB and its external random-number function. A modern implementation should compose typed APIs: sample a sorted combination, then map it through a pure gap transform.}
\cue{Retrieval practice}{\begin{enumerate}\item What is the sampled universe size?\item Why are there $k-1$ separators?\item What is the last sentinel?\item How does telescoping prove the sum?\item Why is the distribution uniform?\end{enumerate}}
\end{longtable}
\begin{tcolorbox}[colback=Ink!4,colframe=Ink,title=Cornell summary]
RANCOM turns the balls-and-bars bijection into a sampler. It uniformly chooses the $k-1$ separator positions among $n+k-1$ available positions, adds endpoints, and records the gap sizes minus one. Strictly increasing separators make every part nonnegative, while telescoping proves the parts total $n$. Since each weak composition has exactly one separator set, uniformity transfers immediately through the bijection. The implementation is compact because the separator array becomes the output array, and the main caution is correct one-based endpoint arithmetic.
\end{tcolorbox}
\end{document}
